
\subsubsection{3.1 Models}

Two models were compared:

\textbf{CodeRL (RL-aligned):} CodeT5-large (770M parameters) with execution-based RL training using binary pass/fail rewards (Salesforce/codet5-large-ntp-py). This model uses an encoder-decoder architecture.

\textbf{CodeLlama-7B-Instruct (instruction-tuned):} CodeLlama with instruction tuning (7B parameters, codellama/CodeLlama-7b-Instruct-hf). This model uses a decoder-only architecture. It is a general instruction-following model, not specifically trained with code-focused DPO.

The models differ in architecture (encoder-decoder vs. decoder-only), parameter count (770M vs. 7B), base pre-training data, and fine-tuning procedure. These differences confound attribution of any observed effects to specific training objectives.

\subsubsection{3.2 Datasets}

Evaluation used HumanEval+ (164 problems) and MBPP+ (378 problems), totaling 542 problems. These benchmarks extend the original HumanEval and MBPP with additional test cases. One sample was generated per problem per model at temperature 0.8, yielding 1,084 total samples.

\subsubsection{3.3 Error Classification}

Errors were classified into three coarse categories based on Python exception types:
\begin{itemize}
\item \textbf{Syntax errors:} Code fails to parse (SyntaxError, IndentationError)
\item \textbf{Runtime errors:} Code parses but raises an exception during execution (TypeError, NameError, ValueError, etc.)
\item \textbf{Assertion errors:} Code executes completely but produces incorrect output (AssertionError from test cases)
\end{itemize}

For fine-grained analysis, the LlmFix 19-category taxonomy was applied, distinguishing specific error causes (e.g., indentation\_error vs. syntax\_error within the syntax category).

\subsubsection{3.4 Execution Depth Measurement}

Execution depth was measured using Python's sys.settrace() function. For each failed sample, the number of unique lines executed before failure was counted and divided by the total number of executable lines in the generated code plus test harness. This provides a measure of how far into the code execution proceeded before failure.

\subsubsection{3.5 Statistical Analysis}

\begin{itemize}
\item \textbf{H-E1 (existence):} Chi-square test for independence on the $2\times 3$ contingency table (model $\times$ error type), with Cramér's V for effect size
\item \textbf{H-M1 (assertion proportion):} Fisher's exact test (one-sided) on the $2\times 2$ table (model $\times$ assertion/non-assertion)
\item \textbf{H-M2 (execution depth):} Welch's t-test (unequal variances assumed) with Cohen's d for effect size
\item \textbf{H-M3 (fine-grained taxonomy):} Chi-square test with Cramér's V at both coarse (3-category) and fine (19-category) granularities
\end{itemize}
